\chapter{Introduction}
\label{ch:intro}

\section{Problem and motivation}
\label{sec:intro-problem}

Software is beginning to be built out of language models of very different capability. A cheap,
small model handles the routine calls; a strong, expensive one is brought in for the hard ones;
and more and more the two are left to coordinate with each other rather than through a person.
When agents of unequal ability share a task with something at stake, a natural assumption rides
along: the capable one will come out ahead. This thesis puts that assumption to a test, and finds
it is not safe.

The test is an old instrument turned on a new subject. Experimental economics has a small set of
games built to isolate the tension between acting for oneself and acting for the group, and to
measure, in tokens, who cooperates and who takes advantage. Put language models in the seats and
the games become a controlled way to ask what a mixed population of agents does. When a strong
model and several weak ones share a public good, does the strong one dominate the weak or
subsidize them? Does letting the agents make public statements before they act help, or hand
someone a lever? And does whatever happens in a group carry over to a one-on-one exchange, or
belong to the crowd?

These are the thesis's three questions, and the short answer it returns to each is that the
intuition is wrong in an instructive way. The capable models do not reliably win; under one
condition they reliably lose. The thing that flips the outcome is not an exchange of messages but
the act of being asked to commit in public. And the gap between what an agent says and what it
does, the nearest thing these agents have to deception, separates the strong from the weak only
when it is measured the right way. What follows is the apparatus that produced those answers and
the case for trusting them.

\section{Research questions and contributions}
\label{sec:intro-rqs}

The study turns on three questions, registered in advance with the tests that would bear on them
and the results read back against those tests (Chapter~\ref{ch:methodology}).

The first asks whether capability asymmetry produces inequality: when strong and weak models share
a game, does welfare spread apart, and does the strong tier earn a premium over the weak
(hypotheses H1 and H2)? The second asks whether communication changes that gap: does giving the
agents a public channel before they act help the weak, the strong, or no one (hypothesis H3)? The
third asks whether any effect is general or local: does what happens in a four-player public good
also happen in a two-player exchange, or reverse between them (hypothesis H5)? Running through all
three is the measurement of deception, the gap between what an agent says and what it does, which
the pre-registration treats as exploratory (hypothesis H4); and a secondary, descriptive question
asks whether an agent's capability tier can be read back from its behavior alone.

The thesis makes two contributions. The first is empirical. The capability premium is real, but its
sign is not fixed: without communication the strong models out-earn the weak, and with it the
premium inverts and the strong earn less. A mixed-effects model localizes that swing to the group
game, and all three strong models show it on their own. The lever is not the exchange of messages
it appears to be, a point developed in Section~\ref{sec:disc-premium}, but the framing of being
asked to commit in public; and the effect belongs to the group game, fading in the two-player
exchange. The second contribution is methodological. Deception measured as a global divergence
between an agent's stated and revealed policies fails in a specific way, scoring the steadiest
cooperators as the least honest; a contemporaneous, action-by-action surprisal does not, and
reporting the two together turns a measurement choice into a result. Neither contribution is the one
the project set out expecting, which is the better reason to trust them: both survived a design
fixed before the data existed.

\section{Scope and limitations}
\label{sec:intro-scope}

Two boundaries belong at the outset. The first is what the communication condition is. It does not
let the agents talk to one another. Each is asked to write a public message and told the others
will see it, but the message is never delivered; agents act on the history of past moves, not on
each other's words (Chapter~\ref{ch:implementation}). The condition is public commitment without
reception, and the claims made about it are claims about that, not about an exchange of cheap talk.
The name is the field's; the apparatus is narrower, and the thesis keeps the two apart throughout.

The second is the modesty of the instrument. ``Capability'' here means three specific frontier
models set against one small local one, not a graded scale; the deception measure is validated only
weakly, against a single human rater, and is reported as exploratory; the whole study is 320 games,
run under a hundred-euro budget on a single consumer GPU. These bounds are stated where they bite
and gathered in Section~\ref{sec:disc-limitations}. They limit how far the numbers travel, not
whether the central effects are real: those are large, consistent across three independent strong
models, and fixed in direction before the first game was played.

\section{Thesis structure}
\label{sec:intro-structure}

The rest of the thesis runs as follows. Chapter~\ref{ch:background} places the work among the
economic games it borrows and the literature on language-model agents around it.
Chapter~\ref{ch:methodology} sets out the design: the populations, the games, the conditions, the
metrics, and the pre-registration that fixed them in advance. Chapter~\ref{ch:implementation}
describes the software that ran it, the agent interface, the game engine, the communication
condition as it actually works, and the measurement pipeline. Chapter~\ref{ch:results} reports what
the 320 games showed, hypothesis by hypothesis, and Chapter~\ref{ch:discussion} interprets it: why
the premium inverts, why only in the group, and what the deception measures do and do not establish.
Chapter~\ref{ch:conclusion} draws the two findings together.
